%% pcp.tex -- short note on PCP(0,poly) vs the PCP theorem
%% %% pcp.tex -- short note on PCP(0,poly) vs the PCP theorem
%% %% pcp.tex -- short note on PCP(0,poly) vs the PCP theorem
%% %% pcp.tex -- short note on PCP(0,poly) vs the PCP theorem
%% \input{pcp} inside zero_knowledge.tex

\subsection*{the PCP theorem and local verification}

\todo{1. revise this and make it fit inside the narative. We need Kilian PCP -> SNARK and maybe some other stuff for NIZK
2. NIZK is impossible without CRS or ROM, need to find the proof
3. what routes SNARG(K)s takes and why(building intuition, details shall be given in its own section)?}

\subsubsection*{Trivial characterization: $\mathrm{NP} = \mathrm{PCP}(0,\mathrm{poly}(n))$}

\textbf{Setting.} Zero randomness, $\mathrm{poly}(n)$ queries.

\begin{itemize}
  \item \textbf{$\mathrm{NP}\subseteq \mathrm{PCP}(0,\mathrm{poly}(n))$:}
  Let proof $\pi := w$ (the NP witness). The verifier has no randomness,
  reads the entire $\pi$ (poly many queries), and simulates the
  deterministic NP-verifier $V(x,\pi)$.
  \begin{itemize}
    \item Completeness: $x\in L \Rightarrow \pi=w$, $V(x,w)=1$, accept w.p.\ 1.
    \item Soundness: $x\notin L \Rightarrow \forall \pi,\ V(x,\pi)=0$, reject w.p.\ 1.
  \end{itemize}

  \item \textbf{$\mathrm{PCP}(0,\mathrm{poly}(n))\subseteq \mathrm{NP}$:}
  No randomness $\Rightarrow$ queried positions are \emph{fixed} (depend only on $x$),
  so only $q(n)=\mathrm{poly}(n)$ positions of $\pi$ ever matter. Take the
  contents of those positions as the NP witness; simulate the deterministic
  verifier on them.
\end{itemize}

This is purely a restatement of the NP definition: \emph{``proof = witness,
read it all, check deterministically.''}

\subsubsection*{The PCP Theorem: $\mathrm{NP} = \mathrm{PCP}(O(\log n), O(1))$}

\textbf{Setting.} $O(\log n)$ random bits, $O(1)$ queries
(constant-query, constant soundness error after amplification).

\begin{itemize}
  \item $O(\log n)$ randomness $\Rightarrow$ proof length $q(n)\cdot 2^{O(\log n)} = \mathrm{poly}(n)$,
        so proofs remain polynomial-size --- consistent with $\mathrm{NP}$.
  \item Despite querying only $O(1)$ symbols, the verifier can detect
        \emph{any} cheating proof with probability $\geq 1/2$ (boostable
        to $1-\varepsilon$ by repetition).
  \item Non-trivial direction: $\mathrm{NP}\subseteq \mathrm{PCP}(O(\log n),O(1))$
        requires constructing an encoding of any NP witness such that
        \emph{local} (constant-query) consistency checks detect \emph{global}
        inconsistency --- the deep result of Arora--Safra and
        Arora--Lund--Motwani--Sudan--Szegedy.
\end{itemize}

\begin{center}
\small
\begin{tabular}{l|c|c}
 & $\mathrm{PCP}(0,\mathrm{poly}(n))$ & $\mathrm{PCP}(O(\log n),O(1))$ \\\hline
Randomness & none & $O(\log n)$ bits \\
Queries & $\mathrm{poly}(n)$ & $O(1)$ \\
Proof length & $\mathrm{poly}(n)$ & $\mathrm{poly}(n)$ \\
Why $= \mathrm{NP}$ & definitional restatement & deep theorem; local $\Rightarrow$ global soundness \\
\end{tabular}
\end{center}

Intuition: the trivial case says \emph{``verify NP by reading the whole witness.''}
The PCP theorem says \emph{``a constant number of cleverly encoded bits suffice,
no matter how large the witness is.''} This is the information-theoretic ancestor
of the IOP-to-SNARK paradigm.
 inside zero_knowledge.tex

\subsection*{the PCP theorem and local verification}

\todo{1. revise this and make it fit inside the narative. We need Kilian PCP -> SNARK and maybe some other stuff for NIZK
2. NIZK is impossible without CRS or ROM, need to find the proof
3. what routes SNARG(K)s takes and why(building intuition, details shall be given in its own section)?}

\subsubsection*{Trivial characterization: $\mathrm{NP} = \mathrm{PCP}(0,\mathrm{poly}(n))$}

\textbf{Setting.} Zero randomness, $\mathrm{poly}(n)$ queries.

\begin{itemize}
  \item \textbf{$\mathrm{NP}\subseteq \mathrm{PCP}(0,\mathrm{poly}(n))$:}
  Let proof $\pi := w$ (the NP witness). The verifier has no randomness,
  reads the entire $\pi$ (poly many queries), and simulates the
  deterministic NP-verifier $V(x,\pi)$.
  \begin{itemize}
    \item Completeness: $x\in L \Rightarrow \pi=w$, $V(x,w)=1$, accept w.p.\ 1.
    \item Soundness: $x\notin L \Rightarrow \forall \pi,\ V(x,\pi)=0$, reject w.p.\ 1.
  \end{itemize}

  \item \textbf{$\mathrm{PCP}(0,\mathrm{poly}(n))\subseteq \mathrm{NP}$:}
  No randomness $\Rightarrow$ queried positions are \emph{fixed} (depend only on $x$),
  so only $q(n)=\mathrm{poly}(n)$ positions of $\pi$ ever matter. Take the
  contents of those positions as the NP witness; simulate the deterministic
  verifier on them.
\end{itemize}

This is purely a restatement of the NP definition: \emph{``proof = witness,
read it all, check deterministically.''}

\subsubsection*{The PCP Theorem: $\mathrm{NP} = \mathrm{PCP}(O(\log n), O(1))$}

\textbf{Setting.} $O(\log n)$ random bits, $O(1)$ queries
(constant-query, constant soundness error after amplification).

\begin{itemize}
  \item $O(\log n)$ randomness $\Rightarrow$ proof length $q(n)\cdot 2^{O(\log n)} = \mathrm{poly}(n)$,
        so proofs remain polynomial-size --- consistent with $\mathrm{NP}$.
  \item Despite querying only $O(1)$ symbols, the verifier can detect
        \emph{any} cheating proof with probability $\geq 1/2$ (boostable
        to $1-\varepsilon$ by repetition).
  \item Non-trivial direction: $\mathrm{NP}\subseteq \mathrm{PCP}(O(\log n),O(1))$
        requires constructing an encoding of any NP witness such that
        \emph{local} (constant-query) consistency checks detect \emph{global}
        inconsistency --- the deep result of Arora--Safra and
        Arora--Lund--Motwani--Sudan--Szegedy.
\end{itemize}

\begin{center}
\small
\begin{tabular}{l|c|c}
 & $\mathrm{PCP}(0,\mathrm{poly}(n))$ & $\mathrm{PCP}(O(\log n),O(1))$ \\\hline
Randomness & none & $O(\log n)$ bits \\
Queries & $\mathrm{poly}(n)$ & $O(1)$ \\
Proof length & $\mathrm{poly}(n)$ & $\mathrm{poly}(n)$ \\
Why $= \mathrm{NP}$ & definitional restatement & deep theorem; local $\Rightarrow$ global soundness \\
\end{tabular}
\end{center}

Intuition: the trivial case says \emph{``verify NP by reading the whole witness.''}
The PCP theorem says \emph{``a constant number of cleverly encoded bits suffice,
no matter how large the witness is.''} This is the information-theoretic ancestor
of the IOP-to-SNARK paradigm.
 inside zero_knowledge.tex

\subsection*{the PCP theorem and local verification}

\todo{1. revise this and make it fit inside the narative. We need Kilian PCP -> SNARK and maybe some other stuff for NIZK
2. NIZK is impossible without CRS or ROM, need to find the proof
3. what routes SNARG(K)s takes and why(building intuition, details shall be given in its own section)?}

\subsubsection*{Trivial characterization: $\mathrm{NP} = \mathrm{PCP}(0,\mathrm{poly}(n))$}

\textbf{Setting.} Zero randomness, $\mathrm{poly}(n)$ queries.

\begin{itemize}
  \item \textbf{$\mathrm{NP}\subseteq \mathrm{PCP}(0,\mathrm{poly}(n))$:}
  Let proof $\pi := w$ (the NP witness). The verifier has no randomness,
  reads the entire $\pi$ (poly many queries), and simulates the
  deterministic NP-verifier $V(x,\pi)$.
  \begin{itemize}
    \item Completeness: $x\in L \Rightarrow \pi=w$, $V(x,w)=1$, accept w.p.\ 1.
    \item Soundness: $x\notin L \Rightarrow \forall \pi,\ V(x,\pi)=0$, reject w.p.\ 1.
  \end{itemize}

  \item \textbf{$\mathrm{PCP}(0,\mathrm{poly}(n))\subseteq \mathrm{NP}$:}
  No randomness $\Rightarrow$ queried positions are \emph{fixed} (depend only on $x$),
  so only $q(n)=\mathrm{poly}(n)$ positions of $\pi$ ever matter. Take the
  contents of those positions as the NP witness; simulate the deterministic
  verifier on them.
\end{itemize}

This is purely a restatement of the NP definition: \emph{``proof = witness,
read it all, check deterministically.''}

\subsubsection*{The PCP Theorem: $\mathrm{NP} = \mathrm{PCP}(O(\log n), O(1))$}

\textbf{Setting.} $O(\log n)$ random bits, $O(1)$ queries
(constant-query, constant soundness error after amplification).

\begin{itemize}
  \item $O(\log n)$ randomness $\Rightarrow$ proof length $q(n)\cdot 2^{O(\log n)} = \mathrm{poly}(n)$,
        so proofs remain polynomial-size --- consistent with $\mathrm{NP}$.
  \item Despite querying only $O(1)$ symbols, the verifier can detect
        \emph{any} cheating proof with probability $\geq 1/2$ (boostable
        to $1-\varepsilon$ by repetition).
  \item Non-trivial direction: $\mathrm{NP}\subseteq \mathrm{PCP}(O(\log n),O(1))$
        requires constructing an encoding of any NP witness such that
        \emph{local} (constant-query) consistency checks detect \emph{global}
        inconsistency --- the deep result of Arora--Safra and
        Arora--Lund--Motwani--Sudan--Szegedy.
\end{itemize}

\begin{center}
\small
\begin{tabular}{l|c|c}
 & $\mathrm{PCP}(0,\mathrm{poly}(n))$ & $\mathrm{PCP}(O(\log n),O(1))$ \\\hline
Randomness & none & $O(\log n)$ bits \\
Queries & $\mathrm{poly}(n)$ & $O(1)$ \\
Proof length & $\mathrm{poly}(n)$ & $\mathrm{poly}(n)$ \\
Why $= \mathrm{NP}$ & definitional restatement & deep theorem; local $\Rightarrow$ global soundness \\
\end{tabular}
\end{center}

Intuition: the trivial case says \emph{``verify NP by reading the whole witness.''}
The PCP theorem says \emph{``a constant number of cleverly encoded bits suffice,
no matter how large the witness is.''} This is the information-theoretic ancestor
of the IOP-to-SNARK paradigm.
 inside zero_knowledge.tex

\subsection*{the PCP theorem and local verification}

\todo{1. revise this and make it fit inside the narative. We need Kilian PCP -> SNARK and maybe some other stuff for NIZK
2. NIZK is impossible without CRS or ROM, need to find the proof
3. what routes SNARG(K)s takes and why(building intuition, details shall be given in its own section)?}

\subsubsection*{Trivial characterization: $\mathrm{NP} = \mathrm{PCP}(0,\mathrm{poly}(n))$}

\textbf{Setting.} Zero randomness, $\mathrm{poly}(n)$ queries.

\begin{itemize}
  \item \textbf{$\mathrm{NP}\subseteq \mathrm{PCP}(0,\mathrm{poly}(n))$:}
  Let proof $\pi := w$ (the NP witness). The verifier has no randomness,
  reads the entire $\pi$ (poly many queries), and simulates the
  deterministic NP-verifier $V(x,\pi)$.
  \begin{itemize}
    \item Completeness: $x\in L \Rightarrow \pi=w$, $V(x,w)=1$, accept w.p.\ 1.
    \item Soundness: $x\notin L \Rightarrow \forall \pi,\ V(x,\pi)=0$, reject w.p.\ 1.
  \end{itemize}

  \item \textbf{$\mathrm{PCP}(0,\mathrm{poly}(n))\subseteq \mathrm{NP}$:}
  No randomness $\Rightarrow$ queried positions are \emph{fixed} (depend only on $x$),
  so only $q(n)=\mathrm{poly}(n)$ positions of $\pi$ ever matter. Take the
  contents of those positions as the NP witness; simulate the deterministic
  verifier on them.
\end{itemize}

This is purely a restatement of the NP definition: \emph{``proof = witness,
read it all, check deterministically.''}

\subsubsection*{The PCP Theorem: $\mathrm{NP} = \mathrm{PCP}(O(\log n), O(1))$}

\textbf{Setting.} $O(\log n)$ random bits, $O(1)$ queries
(constant-query, constant soundness error after amplification).

\begin{itemize}
  \item $O(\log n)$ randomness $\Rightarrow$ proof length $q(n)\cdot 2^{O(\log n)} = \mathrm{poly}(n)$,
        so proofs remain polynomial-size --- consistent with $\mathrm{NP}$.
  \item Despite querying only $O(1)$ symbols, the verifier can detect
        \emph{any} cheating proof with probability $\geq 1/2$ (boostable
        to $1-\varepsilon$ by repetition).
  \item Non-trivial direction: $\mathrm{NP}\subseteq \mathrm{PCP}(O(\log n),O(1))$
        requires constructing an encoding of any NP witness such that
        \emph{local} (constant-query) consistency checks detect \emph{global}
        inconsistency --- the deep result of Arora--Safra and
        Arora--Lund--Motwani--Sudan--Szegedy.
\end{itemize}

\begin{center}
\small
\begin{tabular}{l|c|c}
 & $\mathrm{PCP}(0,\mathrm{poly}(n))$ & $\mathrm{PCP}(O(\log n),O(1))$ \\\hline
Randomness & none & $O(\log n)$ bits \\
Queries & $\mathrm{poly}(n)$ & $O(1)$ \\
Proof length & $\mathrm{poly}(n)$ & $\mathrm{poly}(n)$ \\
Why $= \mathrm{NP}$ & definitional restatement & deep theorem; local $\Rightarrow$ global soundness \\
\end{tabular}
\end{center}

Intuition: the trivial case says \emph{``verify NP by reading the whole witness.''}
The PCP theorem says \emph{``a constant number of cleverly encoded bits suffice,
no matter how large the witness is.''} This is the information-theoretic ancestor
of the IOP-to-SNARK paradigm.
